\section{Beschreibung des Unternehmens}
\label{sec:unternehmen-rdy}

Dieser Abschnitt berichtet. Er bewertet nicht. Alle Betr\"age in Mio.\,INR,
sofern nicht anders bezeichnet; Betr\"age je Aktie in INR. Das
Gesch\"aftsjahr endet am 31.~M\"arz, FY2026 meint das Jahr zum 31.~M\"arz 2026.

\subsection{Handelnde Personen}

Vorsitzender des Board ist K.~Satish Reddy, seit 1993 im Unternehmen.
Co-Vorsitzender und Managing Director ist G.~V.~Prasad, seit 1990 im
Unternehmen und im Juli 2025 f\"ur f\"unf weitere Jahre
wiederbestellt.\quelleZF{2026}{63}\quelleMerken{s23zfxcacgxgd}
Chief Executive Officer ist seit 2018 Erez Israeli, Chief Financial Officer
seit 2000 M.~V.~Narasimham.\quelleErneut{s23zfxcacgxgd}
Dem Board geh\"oren \Stueck{\RdyDirektoren} Direktoren an, davon
\Stueck{\RdyDirektorenUnabh} nach den Regeln der New York Stock Exchange
unabh\"angig; es unterh\"alt \Stueck{\RdyAusschuesse}
Aussch\"usse.\quelleZF{2026}{72}\quelleMerken{s23zfxcacgxhc}\quelleZF{2026}{73}\quelleMerken{s23zfxcacgxhd}

Eine Tatsache geh\"ort hierher, weil der Bericht sie selbst ausweist und weil
sie die Verflechtung von Familie und Organ beschreibt: Der Vorsitzende und der
Co-Vorsitzende sind Schw\"ager; der Bericht vermerkt es wechselseitig in den
Fu{\ss}noten der Direktorentabelle.\quelleErneut{s23zfxcacgxgd}
Beide sitzen zugleich im Nachhaltigkeits- und im
Aktion\"arsausschuss.\quelleZF{2026}{76}\quelleMerken{s23zfxcacgxhg}
Die Ehefrauen und die Kinder des Co-Vorsitzenden sind im Bericht als
\enquote{Family Members} der Promoter-Gruppe gef\"uhrt.\quelleZF{2026}{80}\quelleMerken{s23zfxcacgxia}

Die Verg\"utung der beiden Vollzeitdirektoren betrug \MioINR{\RdyVerguetungChairman}
und \MioINR{\RdyVerguetungCoChairman}; der Chief Executive Officer erhielt
\MioINR{\RdyVerguetungCeo} in bar, der Chief Financial Officer
\MioINR{\RdyVerguetungCfo}.\quelleZF{2026}{71}\quelleMerken{s23zfxcacgxhb}
Die Kommission der Vollzeitdirektoren ist gesetzlich auf je
\Pct{\RdyDeckelVollzeit} des Jahresnettogewinns
begrenzt.\quelleZF{2026}{70}\quelleMerken{s23zfxcacgxha}

\subsection{Eigent\"umerstruktur}

Zum 31.~M\"arz 2026 waren \Stueck{\RdyAktienAusgegeben} Aktien mit einem
Nennwert von je einer Rupie ausgegeben.\quelleZF{2026}{81}\quelleMerken{s23zfxcacgxib}

\begin{table}[htbp]
\centering
\caption{Eigent\"umerstruktur zum 31.~M\"arz 2026}
\label{tab:rdy-eigentuemer}
\begin{tabular}{lr}
\toprule
Gruppe & Anteil\\
\midrule
Promoter-Gruppe (Gr\"underfamilien)        & \Pct{\RdyPromoterAnteil}\\
Life Insurance Corporation of India        & \Pct{\RdyLicAnteil}\\
\"ubriger Streubesitz                       & \Pct{\RdyStreubesitzOhneLic}\\
\midrule
Summe                                       & \Pct{\RdySumme}\\
\bottomrule
\end{tabular}
\end{table}

Anders als bei vielen Emittenten geht diese Rechnung auf: Das Unternehmen
weist Promoter-Gruppe und Streubesitz selbst aus, und beide summieren sich zu
hundert Prozent.\quelleErneut{s23zfxcacgxib}
Als Hinterlegungsscheine gehalten werden \Pct{\RdyAdrAnteil} der Aktien; ein
Schein entspricht einer Aktie, und dieses Verh\"altnis hat sich durch den
Aktiensplit nicht ge\"andert.\quelleZF{2026}{82}\quelleMerken{s23zfxcacgxic}
In Indien sind \Stueck{\RdyRecordHolder} Inhaber eingetragen.\quelleErneut{s23zfxcacgxic}

\subsection{Beurteilung durch die Kapitalm\"arkte}

Die Aktie ist an der BSE und der National Stock Exchange of India notiert, der
Hinterlegungsschein seit 2001 an der New York Stock Exchange unter dem
K\"urzel RDY.\quelleZF{2026}{83}\quelleMerken{s24zfxcacgxid}
Zum 28.~Oktober 2024 wurde die Aktie im Verh\"altnis eins zu f\"unf geteilt;
der Nennwert sank von f\"unf Rupien auf eine, die Aktienzahl stieg von
\Stueck{\RdyAktienVorSplit} auf \Stueck{\RdyAktienNachSplit}.\quelleZF{2026}{157}\quelleMerken{s24zfxcacgxbfh}
Alle Angaben je Aktie in diesem Bericht sind auf die Zeit nach dem Split
umgerechnet; das Unternehmen tut dies in seinem Abschluss
ebenso.\quelleErneut{s24zfxcacgxbfh}

Bei einem Schlusskurs des Hinterlegungsscheins von \USDje{\RdyKurs} am
\Datum{\RdyKursTag} entspricht das \INRje{\RdyKursRupien} je Aktie und einem
B\"orsenwert von \MrdINR{\RdyMarktwertMrd}. Das Kurs-Gewinn-Verh\"altnis
betr\"agt \Faktor{\RdyKursGewinn}, das Kurs-Buchwert-Verh\"altnis
\Faktor{\RdyKursBuchwert}.

Eine Dividendenpolitik mit einer Quote oder einem Ziel gibt es nicht. Der
Bericht formuliert nur ein Verfahren: \enquote{Every year our Board of
Directors recommends the amount of dividends to be paid to shareholders, if
any, based upon conditions then existing.}\quelleZF{2026}{82}\quelleMerken{s24zfxcacgxic}
Gezahlt wurden zuletzt \INRje{\RdyDividendeJeAktie} je Aktie, derselbe
Betrag wie in den beiden Vorjahren, nach dem Split
gerechnet.\quelleErneut{s24zfxcacgxic}
Das entspricht einer Rendite von \Pct{\RdyDividendenrendite} und einer
Aussch\"uttungsquote von \Pct{\RdyAusschuettungsquote}.

\subsection{Analystenabdeckung und Konsens}

\Stueck{\RdyKonsensHaeuser} H\"auser bilden den von S\&P Global gesammelten
Konsens f\"ur den Hinterlegungsschein. Er lautet \enquote{Hold} bei einem
mittleren Kursziel von \USDje{\RdyKonsensZiel}, \Pct{\RdyKonsensAufschlag}
\"uber dem Bezugskurs.\quelleWeb{stockanalysis.com, Analystenkonsens RDY (Datenbasis S\&P Global)}{quelle:konsensrdy}{18.09.2026}\quelleMerken{s24webxstockanalysisxcomxxanalystenkonsensxrdyxxdatenbasisxsxxpxglobalxxquellexkonsensrdy}
Auch diese Angabe besteht die Probe gegen sich selbst: Aus Kursziel und
Aufschlag folgt ein Bezugskurs von \USDje{\RdyKonsensBezug}, und das ist der
Schlusskurs des \Datum{\RdyKursTag}.

Die Abdeckung des Hinterlegungsscheins ist erkennbar d\"unn. Die indische
Notierung wird von mehr H\"ausern verfolgt; deren Sch\"atzungen sind hier
nicht erhoben, weil sie sich auf ein anderes Papier und eine andere
W\"ahrung beziehen. Dass f\"unf Sch\"atzungen eine schw\"achere Grundlage
sind als elf, geh\"ort zum Befund und wird in
Abschnitt~\ref{sec:verdikt-rdy} nicht \"ubergangen.

\subsection{Nicht verf\"ugbare Angaben}

Zwei Angaben fehlen, und die zweite ist die wichtigste L\"ucke dieses ganzen
Dokuments.

Eine ISIN oder CUSIP nennt keine der ausgewerteten Quellen. Das ist
belanglos.

\emph{Der Umsatzbeitrag des Lenalidomid-Gesch\"afts ist nirgends beziffert.}
Gesucht wurde in f\"unf Einreichungen \"uber sieben Gesch\"aftsjahre nach den
Begriffen \enquote{Lenalidomide} und \enquote{Revlimid}. Das Unternehmen
nennt eine Preissenkung und ihre Wirkung auf einen einzelnen Posten
(Abschnitt~\ref{sec:strategie-rdy}), aber weder einen Umsatz noch einen
Ergebnisbeitrag noch einen Zeitpunkt des Auslaufens. Die einzige zeitliche
Angabe stammt aus der Wiedergabe einer gegnerischen Klageschrift, nicht aus
einer eigenen Aussage des Unternehmens.\quelleZF{2026}{187}\quelleMerken{s25zfxcacgxbih}
Das ist keine unvollst\"andige Suche, sondern eine Angabe, die es nicht gibt,
und sie ist genau die Gr\"o{\ss}e, an der die Bewertung dieses Titels
h\"angt.
